
On
obtient encore essentiellement les mêmes résultats pour \(\Gr^\bullet\) au lieu de \(K_\bullet\) ou \(K^\bullet\), avec la même réserve concernant l'assertion de bijectivité dans (i). Dans le cas où on se borne à \(\Gr^1=\Pic\), cette assertion de bijectivité est valable si \(X\) est noethérien et \emph{normal}, car dans ce cas, pour tout ouvert \(U\) d'un espace affine sur \(k\), on voit aussitôt que \(\Pic(X)\longrightarrow \Pic(X\times_k U)\) est surjectif (EGA \(\mathrm{Err}_{\mathrm{IV}}\) 53, 21.4.13).

\subsection*{BIBLIOGRAPHIE}
\begin{enumerate}[label={[\arabic*]}]
\item Dickson, E. Linear groups, Dover, New York, 1958.
\item Grothendieck, A. TDTE, V, VI~; Séminaire Bourbaki 232, 236.
\item Kleiman, S. Towards a numerical theory of ampleness, Ann. of Math. 84 (1966), pp. 293-344.
\item Murre, J. On contravariant functors from the category of préschemes over a field into the category of abelian groups, Publications Mathématiques de l'IHES, 23, 1964.
\item Chevalley, C. Les classes d'équivalence rationnelles, I et II, Séminaire à l'ENS 1958, exposés 2 et 3.
\item S. Abhyankar, Resolution of singularities of arithmetical surfaces, in~: Arithmetical Algebraic Geometry, Purdue University, Harper--Row, New York 1965.
\end{enumerate}

\providecommand{\Isom}{\operatorname{Isom}}
\providecommand{\Hom}{\operatorname{Hom}}
\providecommand{\Coker}{\operatorname{Coker}}
\providecommand{\Nil}{\operatorname{Nil}}
\providecommand{\im}{\operatorname{Im}}
\providecommand{\coker}{\operatorname{coker}}
\providecommand{\Supp}{\operatorname{Supp}}
\providecommand{\fppf}{\mathrm{fppf}}
\providecommand{\zar}{\mathrm{zar}}
\providecommand{\GL}{\operatorname{Gl}}
\providecommand{\Gr}{\operatorname{Gr}}
\providecommand{\Pic}{\operatorname{Pic}}
\providecommand{\Spec}{\operatorname{Spec}}
\providecommand{\V}{\operatorname{V}}
\providecommand{\Ass}{\operatorname{Ass}}

\newpage
